\documentclass[11pt]{article}

\usepackage[margin=1in]{geometry}
\usepackage{amsmath, amssymb, amsfonts, bm, mathtools}
\usepackage{enumitem}

\newcommand{\Rbar}{\bar R}
\newcommand{\nuo}{\nu_0}
\newcommand{\sigb}{\sigma_0^2}
\newcommand{\E}{\mathbb E}
\newcommand{\KL}{\operatorname{KL}}
\newcommand{\Normal}{\mathcal N}
\newcommand{\IW}{\operatorname{IW}}
\newcommand{\GIG}{\operatorname{GIG}}
\newcommand{\chisq}{\chi^2}

\title{Collapsed Likelihood and VEB Algorithm for RSS Fine-Mapping with LD Uncertainty}
\author{}
\date{}

\begin{document}
\maketitle

\section{Original model}

Let \(v \in \mathbb R^p\) denote the vector of summary statistics,
for example
\[
v = \frac{1}{N}X^\top y,
\]
and let \(\Rbar\) be an out-of-sample reference LD matrix.

We consider the model
\begin{align}
R &\sim \IW_p(\nuo \Rbar,\nuo+p+1), \label{eq:prior_R}\\
v \mid R,b &\sim \Normal_p(Rb,R/N). \label{eq:rss_model}
\end{align}

The inverse-Wishart parameterization is chosen so that
\[
\E(R)=\Rbar.
\]
Indeed, if
\[
R \sim \IW_p(S,m),
\]
then
\[
\E(R)=\frac{S}{m-p-1},
\]
and here
\[
S=\nuo\Rbar,
\qquad
m=\nuo+p+1.
\]

The likelihood is
\[
p(v\mid R,b)
=
\left(\frac{N}{2\pi}\right)^{p/2}
|R|^{-1/2}
\exp\left\{
-\frac N2(v-Rb)^\top R^{-1}(v-Rb)
\right\}.
\]

The quadratic term expands as
\begin{align}
(v-Rb)^\top R^{-1}(v-Rb)
&=
v^\top R^{-1}v
-2b^\top v
+b^\top Rb.
\end{align}

Therefore,
\begin{align}
p(v\mid R,b)
&=
p(v\mid R,b=0)
\exp\left\{
Nb^\top v-\frac N2 b^\top Rb
\right\}. \label{eq:change_measure_identity}
\end{align}

\section{Collapsed likelihood}

The marginal likelihood is
\[
p(v\mid b,\nuo)
=
\int p(v\mid R,b)p(R\mid \nuo)\,dR.
\]

Using \eqref{eq:change_measure_identity},
\begin{align}
p(v\mid b,\nuo)
&=
\exp(Nb^\top v)
\int
p(v\mid R,b=0)
\exp\left\{
-\frac N2 b^\top Rb
\right\}
p(R\mid \nuo)\,dR.
\end{align}

Define the null marginal likelihood
\[
p_0(v\mid \nuo)
=
\int p(v\mid R,b=0)p(R\mid \nuo)\,dR.
\]

Then
\begin{align}
p(v\mid b,\nuo)
&=
p_0(v\mid \nuo)
\exp(Nb^\top v)
\E_{R\mid v,b=0,\nuo}
\left[
\exp\left\{
-\frac N2 b^\top Rb
\right\}
\right]. \label{eq:collapsed_R_expectation}
\end{align}

Under the null model,
\[
v\mid R,b=0 \sim \Normal_p(0,R/N),
\]
and by inverse-Wishart conjugacy,
\[
R\mid v,b=0,\nuo
\sim
\IW_p\left(
\nuo\Rbar+Nvv^\top,
\nuo+p+2
\right).
\]

Let
\[
B_\nu=\nuo\Rbar+Nvv^\top.
\]
If
\[
R\sim \IW_p(B_\nu,\nuo+p+2),
\]
then for fixed \(b\),
\[
b^\top Rb
\overset d=
\frac{b^\top B_\nu b}{\omega},
\qquad
\omega\sim \chisq_{\nuo+3}.
\]

Hence the exact collapsed likelihood is
\begin{align}
p(v\mid b,\nuo)
&=
p_0(v\mid \nuo)
\exp(Nb^\top v)
\E_{\omega\sim \chisq_{\nuo+3}}
\left[
\exp\left\{
-\frac{N}{2\omega}
b^\top(\nuo\Rbar+Nvv^\top)b
\right\}
\right].
\label{eq:collapsed_likelihood}
\end{align}

Define
\[
Q_{\nuo}(b)
=
b^\top(\nuo\Rbar+Nvv^\top)b.
\]
Then
\[
p(v\mid b,\nuo)
=
p_0(v\mid \nuo)
\exp(Nb^\top v)
\E_{\omega\sim \chisq_{\nuo+3}}
\left[
\exp\left\{
-\frac{NQ_{\nuo}(b)}{2\omega}
\right\}
\right].
\]

The null marginal likelihood is a multivariate \(t\) density:
\[
p_0(v\mid \nuo)
=
t_{\nuo+2}
\left(
v;
0,
\frac{\nuo}{N(\nuo+2)}\Rbar
\right).
\]

Equivalently, if
\[
r_0 = v^\top \Rbar^{-1}v,
\qquad
\Sigma_\nu=\frac{\nuo}{N(\nuo+2)}\Rbar,
\]
then
\begin{align}
\log p_0(v\mid \nuo)
&=
\log\Gamma\left(\frac{\nuo+p+2}{2}\right)
-
\log\Gamma\left(\frac{\nuo+2}{2}\right)
-
\frac p2\log\{(\nuo+2)\pi\}
-
\frac12\log|\Sigma_\nu|
\nonumber\\
&\quad
-
\frac{\nuo+p+2}{2}
\log\left(
1+\frac N{\nuo}r_0
\right).
\label{eq:log_p0}
\end{align}

\section{Variational lower bound}

For any variational density \(q(\omega)\),
\begin{align}
\log
\E_{\omega\sim \chisq_{\nuo+3}}
\left[
\exp\left\{
-\frac{NQ_{\nuo}(b)}{2\omega}
\right\}
\right]
&\geq
-\frac N2
\E_q(\omega^{-1})
Q_{\nuo}(b)
-
\KL\{q(\omega)\|\chisq_{\nuo+3}\}.
\end{align}

Therefore,
\begin{align}
\log p(v\mid b,\nuo)
&\geq
\log p_0(v\mid \nuo)
+
Nb^\top v
-
\frac N2
\E_q(\omega^{-1})
Q_{\nuo}(b)
-
\KL\{q(\omega)\|\chisq_{\nuo+3}\}.
\label{eq:basic_elbo}
\end{align}

This bound becomes tight for the optimal density
\[
q^\star(\omega)
\propto
\chisq_{\nuo+3}(\omega)
\exp\left\{
-\frac{NQ_{\nuo}(b)}{2\omega}
\right\}.
\]

Since
\[
\chisq_d(\omega)
\propto
\omega^{d/2-1}\exp(-\omega/2),
\]
we have
\[
q^\star(\omega)
=
\GIG\left(
\lambda=\frac{\nuo+3}{2},
\chi=NQ_{\nuo}(b),
\psi=1
\right),
\]
where the GIG density is
\[
\GIG(\lambda,\chi,\psi)
\propto
\omega^{\lambda-1}
\exp\left\{
-\frac12
\left(
\frac{\chi}{\omega}+\psi\omega
\right)
\right\}.
\]

\section{Single-effect model}

Now consider \(L=1\). Let
\[
\gamma\in\{1,\dots,p\},
\qquad
b=\beta e_\gamma.
\]

Assume
\[
P(\gamma=j)=\pi_j,
\qquad
\sum_{j=1}^p \pi_j=1,
\]
and
\[
\beta\mid \gamma=j,\sigb \sim \Normal(0,\sigb).
\]

For \(\gamma=j\),
\[
b=\beta e_j,
\]
so
\[
b^\top v=\beta v_j.
\]

Also,
\begin{align}
Q_{\nuo}(\beta e_j)
&=
\beta^2 e_j^\top(\nuo\Rbar+Nvv^\top)e_j \\
&=
\beta^2
\left(
\nuo \Rbar_{jj}+Nv_j^2
\right).
\end{align}

Define
\[
a_j(\nuo)=\nuo\Rbar_{jj}+Nv_j^2.
\]
If \(\Rbar\) is a correlation matrix, then typically \(\Rbar_{jj}=1\), giving
\[
a_j(\nuo)=\nuo+Nv_j^2.
\]

The exact collapsed likelihood conditional on \((\gamma=j,\beta)\) is
\begin{align}
p(v\mid \gamma=j,\beta,\nuo)
&=
p_0(v\mid\nuo)
\exp(N\beta v_j)
\E_{\omega\sim\chisq_{\nuo+3}}
\left[
\exp\left\{
-\frac{N a_j(\nuo)\beta^2}{2\omega}
\right\}
\right].
\end{align}

The augmented joint density is
\begin{align}
p(v,\gamma=j,\beta,\omega\mid \nuo,\sigb)
&=
p_0(v\mid\nuo)
\pi_j
\Normal(\beta;0,\sigb)
\chisq_{\nuo+3}(\omega)
\nonumber\\
&\quad\times
\exp\left\{
N\beta v_j
-
\frac{N a_j(\nuo)\beta^2}{2\omega}
\right\}.
\label{eq:augmented_joint_L1}
\end{align}

\section{Mean-field variational family}

Use the variational family
\[
q(\gamma=j,\beta,\omega)
=
\alpha_j q_j(\beta)q_j(\omega),
\]
where
\[
\alpha_j\geq 0,
\qquad
\sum_{j=1}^p \alpha_j=1.
\]

This should be read as
\[
q(\gamma=j)=\alpha_j,
\qquad
q(\beta,\omega\mid \gamma=j)=q_j(\beta)q_j(\omega).
\]

The factorization allows the variational distributions of \(\beta\) and \(\omega\) to depend on the candidate causal SNP \(j\). This is important because each \(j\) has a different \(v_j\) and \(a_j(\nuo)\).

Take
\[
q_j(\beta)=\Normal(\mu_j,s_j^2),
\]
and
\[
q_j(\omega)
=
\GIG(\lambda_j,\chi_j,\psi_j).
\]

In the coordinate-optimal update below,
\[
\lambda_j=\frac{\nuo+3}{2},
\qquad
\psi_j=1.
\]

Define
\[
m_{2j}
=
\E_{q_j}(\beta^2)
=
\mu_j^2+s_j^2,
\]
and
\[
\eta_j
=
\E_{q_j}(\omega^{-1}).
\]

The ELBO is
\begin{align}
\mathcal L(\nuo,\sigb,q)
&=
\log p_0(v\mid\nuo)
+
\sum_{j=1}^p
\alpha_j
\left[
\log \pi_j
-\log \alpha_j
+
\mathcal F_j
\right],
\label{eq:global_elbo}
\end{align}
where
\begin{align}
\mathcal F_j
&=
\E_{q_j(\beta)}
\left[
\log \Normal(\beta;0,\sigb)
-
\log q_j(\beta)
\right]
+
Nv_j\E_{q_j}(\beta)
\nonumber\\
&\quad
-
\frac N2
a_j(\nuo)
\E_{q_j}(\beta^2)
\E_{q_j}(\omega^{-1})
-
\KL\{q_j(\omega)\|\chisq_{\nuo+3}\}.
\label{eq:Fj_general}
\end{align}

Since \(q_j(\beta)=\Normal(\mu_j,s_j^2)\),
\begin{align}
\E_{q_j(\beta)}
\left[
\log \Normal(\beta;0,\sigb)
-
\log q_j(\beta)
\right]
&=
\frac12
\left[
\log\left(\frac{s_j^2}{\sigb}\right)
+
1
-
\frac{m_{2j}}{\sigb}
\right].
\end{align}

Therefore,
\begin{align}
\mathcal F_j
&=
\frac12
\left[
\log\left(\frac{s_j^2}{\sigb}\right)
+
1
-
\frac{m_{2j}}{\sigb}
\right]
+
Nv_j\mu_j
-
\frac N2 a_j(\nuo)m_{2j}\eta_j
-
\KL\{q_j(\omega)\|\chisq_{\nuo+3}\}.
\label{eq:Fj_normal}
\end{align}

\section{Coordinate updates for fixed \(\nuo\) and \(\sigb\)}

For fixed \(\nuo\), \(\sigb\), and \(q_j(\omega)\), the optimal \(q_j(\beta)\) satisfies
\[
q_j(\beta)
\propto
\Normal(\beta;0,\sigb)
\exp\left\{
Nv_j\beta
-
\frac N2 a_j(\nuo)\eta_j\beta^2
\right\}.
\]

Thus
\[
q_j(\beta)=\Normal(\mu_j,s_j^2),
\]
where
\begin{align}
s_j^2
&=
\left[
\frac{1}{\sigb}
+
N a_j(\nuo)\eta_j
\right]^{-1},
\label{eq:s_update}
\\
\mu_j
&=
s_j^2 Nv_j.
\label{eq:mu_update}
\end{align}

Then
\[
m_{2j}=\mu_j^2+s_j^2.
\]

For fixed \(q_j(\beta)\), the optimal \(q_j(\omega)\) is
\begin{align}
q_j(\omega)
&\propto
\chisq_{\nuo+3}(\omega)
\exp\left\{
-\frac{N a_j(\nuo)m_{2j}}{2\omega}
\right\}.
\end{align}

Therefore,
\[
q_j(\omega)
=
\GIG\left(
\lambda_j,
\chi_j,
\psi_j
\right),
\]
with
\begin{align}
\lambda_j
&=
\frac{\nuo+3}{2},
\\
\chi_j
&=
N a_j(\nuo)m_{2j},
\\
\psi_j
&=
1.
\end{align}

For this GIG distribution,
\[
\eta_j
=
\E_{q_j}(\omega^{-1})
=
\left(\frac{\psi_j}{\chi_j}\right)^{1/2}
\frac{
K_{\lambda_j-1}(\sqrt{\chi_j\psi_j})
}{
K_{\lambda_j}(\sqrt{\chi_j\psi_j})
}.
\]

Since \(\psi_j=1\),
\[
\eta_j
=
\frac{1}{\sqrt{\chi_j}}
\frac{
K_{(\nuo+1)/2}(\sqrt{\chi_j})
}{
K_{(\nuo+3)/2}(\sqrt{\chi_j})
}.
\]

If \(\chi_j=0\), use the limit
\[
\eta_j=\frac{1}{\nuo+1}.
\]

For fixed local factors, the optimal posterior inclusion weights are
\[
\alpha_j
=
\frac{
\pi_j\exp(\mathcal F_j)
}{
\sum_{k=1}^p \pi_k\exp(\mathcal F_k)
}.
\]

The variational posterior for \(b\) is therefore
\[
q(b)
=
\sum_{j=1}^p
\alpha_j
\int
\delta_{\beta e_j}(b)
q_j(\beta)
\,d\beta.
\]

Consequently,
\[
\E_q(b_j)=\alpha_j\mu_j,
\]
and
\[
\E_q(b_j^2)=\alpha_j m_{2j}.
\]

\section{KL term for \(q_j(\omega)\)}

Let
\[
d=\nuo+3,
\qquad
\lambda=\frac d2=\frac{\nuo+3}{2}.
\]

The chi-square density is
\[
\chisq_d(\omega)
=
\frac{1}{2^\lambda\Gamma(\lambda)}
\omega^{\lambda-1}
\exp(-\omega/2).
\]

The coordinate-optimal variational density is
\[
q_j(\omega)
=
\GIG(\lambda,\chi_j,1).
\]

Its density is
\[
q_j(\omega)
=
\frac{
\chi_j^{-\lambda/2}
}{
2K_\lambda(\sqrt{\chi_j})
}
\omega^{\lambda-1}
\exp\left\{
-\frac12
\left(
\frac{\chi_j}{\omega}+\omega
\right)
\right\}.
\]

Therefore,
\begin{align}
\KL\{q_j(\omega)\|\chisq_d\}
&=
\log\Gamma(\lambda)
+
(\lambda-1)\log 2
-
\frac{\lambda}{2}\log\chi_j
-
\log K_\lambda(\sqrt{\chi_j})
-
\frac{\chi_j}{2}\eta_j.
\label{eq:KL_omega}
\end{align}

When \(\chi_j\to 0\), this KL goes to zero, as expected.

\section{Empirical-Bayes updates for \(\nuo\) and \(\sigb\)}

We treat \(\nuo\) and \(\sigb\) as empirical-Bayes hyperparameters.

The variational empirical-Bayes objective is
\[
\widehat\ell_{\mathrm{VEB}}(\nuo,\sigb)
=
\max_q \mathcal L(\nuo,\sigb,q).
\]

For fixed \(\nuo\), the update for \(\sigb\) has a closed form. Since
\[
\sum_{j=1}^p \alpha_j=1,
\]
the part of the ELBO depending on \(\sigb\) is
\[
-\frac12
\sum_{j=1}^p
\alpha_j
\left[
\log \sigb+\frac{m_{2j}}{\sigb}
\right].
\]

Maximizing with respect to \(\sigb\) gives
\[
\boxed{
\sigb
\leftarrow
\sum_{j=1}^p
\alpha_j m_{2j}.
}
\label{eq:sigma_update}
\]

There is generally no closed-form update for \(\nuo\), because \(\nuo\) appears in
\[
\log p_0(v\mid\nuo),
\]
\[
a_j(\nuo)=\nuo\Rbar_{jj}+Nv_j^2,
\]
\[
\lambda_j=\frac{\nuo+3}{2},
\]
and
\[
\KL\{q_j(\omega)\|\chisq_{\nuo+3}\}.
\]

Thus update \(\nuo\) numerically:
\[
\boxed{
\nuo
\leftarrow
\arg\max_{\nu>0}
\mathcal L(\nu,\sigb,q).
}
\]

In practice, it is often more stable to optimize over
\[
\rho=\log\nuo,
\]
so that
\[
\nuo=\exp(\rho).
\]

Alternatively, use a one-dimensional grid:
\[
\nuo\in\{\nu^{(1)},\dots,\nu^{(G)}\},
\]
and choose the grid point with the largest optimized ELBO.

\section{VEB algorithm for \(L=1\)}

\begin{enumerate}[leftmargin=2em]

\item \textbf{Input:}
summary statistic \(v\), sample size \(N\), reference LD matrix \(\Rbar\), prior weights \(\pi_j\), initial values \(\nuo^{(0)}\) and \({\sigb}^{(0)}\).

\item \textbf{Initialize:}
for each \(j\), initialize
\[
\alpha_j^{(0)}=\pi_j,
\qquad
\eta_j^{(0)}=\frac{1}{\nuo^{(0)}+1}.
\]

\item \textbf{Repeat until convergence:}

\begin{enumerate}[leftmargin=2em]

\item For each SNP \(j=1,\dots,p\), compute
\[
a_j(\nuo)=\nuo\Rbar_{jj}+Nv_j^2.
\]

\item Update \(q_j(\beta)\):
\[
s_j^2
=
\left[
\frac{1}{\sigb}
+
N a_j(\nuo)\eta_j
\right]^{-1},
\]
\[
\mu_j=s_j^2Nv_j,
\]
\[
m_{2j}=\mu_j^2+s_j^2.
\]

\item Update \(q_j(\omega)\):
\[
\lambda_j=\frac{\nuo+3}{2},
\qquad
\chi_j=N a_j(\nuo)m_{2j},
\qquad
\psi_j=1.
\]
Set
\[
q_j(\omega)=\GIG(\lambda_j,\chi_j,1).
\]

\item Update
\[
\eta_j
=
\E_{q_j}(\omega^{-1})
=
\frac{1}{\sqrt{\chi_j}}
\frac{
K_{(\nuo+1)/2}(\sqrt{\chi_j})
}{
K_{(\nuo+3)/2}(\sqrt{\chi_j})
}.
\]
If \(\chi_j\) is numerically close to zero, use
\[
\eta_j=\frac{1}{\nuo+1}.
\]

\item Compute the local objective
\[
\mathcal F_j
=
\frac12
\left[
\log\left(\frac{s_j^2}{\sigb}\right)
+
1
-
\frac{m_{2j}}{\sigb}
\right]
+
Nv_j\mu_j
-
\frac N2a_j(\nuo)m_{2j}\eta_j
-
\KL\{q_j(\omega)\|\chisq_{\nuo+3}\}.
\]

\item Update posterior inclusion probabilities:
\[
\alpha_j
=
\frac{
\pi_j\exp(\mathcal F_j)
}{
\sum_{k=1}^p
\pi_k\exp(\mathcal F_k)
}.
\]

\item Update the prior variance:
\[
\sigb
\leftarrow
\sum_{j=1}^p
\alpha_j m_{2j}.
\]

\item Update the LD-mismatch parameter:
\[
\nuo
\leftarrow
\arg\max_{\nu>0}
\mathcal L(\nu,\sigb,q).
\]
This step is one-dimensional and can be done by grid search or numerical optimization over \(\log\nu\).

\end{enumerate}

\item \textbf{Output:}
\[
\widehat\nuo,
\qquad
\widehat{\sigb},
\qquad
\alpha_j,
\qquad
q_j(\beta)=\Normal(\mu_j,s_j^2),
\qquad
q_j(\omega)=\GIG(\lambda_j,\chi_j,1).
\]

The posterior mean of the effect vector is
\[
\E_q(b_j)=\alpha_j\mu_j.
\]

The posterior second moment is
\[
\E_q(b_j^2)=\alpha_j(\mu_j^2+s_j^2).
\]

\end{enumerate}

\section{Variant with a shared \(q(\omega)\)}

A simpler approximation uses a single variational distribution for \(\omega\),
shared across all candidate causal SNPs:
\[
q(\gamma=j,\beta,\omega)
=
\alpha_j q_j(\beta)q(\omega).
\]
Write
\[
q(\omega)=\GIG(\lambda,\chi,\psi),
\qquad
\eta=\E_q(\omega^{-1}).
\]

The ELBO becomes
\begin{align}
\mathcal L_{\mathrm{shared}}(\nuo,\sigb,q)
&=
\log p_0(v\mid\nuo)
-
\KL\{q(\omega)\|\chisq_{\nuo+3}\}
\nonumber\\
&\quad+
\sum_{j=1}^p
\alpha_j
\left[
\log \pi_j
-
\log \alpha_j
+
\mathcal G_j
\right],
\end{align}
where
\begin{align}
\mathcal G_j
&=
\frac12
\left[
\log\left(\frac{s_j^2}{\sigb}\right)
+
1
-
\frac{m_{2j}}{\sigb}
\right]
+
Nv_j\mu_j
-
\frac N2 a_j(\nuo)m_{2j}\eta.
\end{align}
The only difference from \eqref{eq:Fj_normal} is that the same \(\eta\) is used for every \(j\), and the KL penalty for \(\omega\) is global rather than local.

For fixed \(q(\omega)\), the update for \(q_j(\beta)\) is
\begin{align}
s_j^2
&=
\left[
\frac{1}{\sigb}
+
N a_j(\nuo)\eta
\right]^{-1},
\\
\mu_j
&=
s_j^2Nv_j,
\\
m_{2j}
&=
\mu_j^2+s_j^2.
\end{align}

For fixed \(\{\alpha_j,q_j(\beta)\}_{j=1}^p\), the coordinate-optimal shared \(q(\omega)\) is
\[
q(\omega)
\propto
\chisq_{\nuo+3}(\omega)
\exp\left\{
-
\frac{N}{2\omega}
\sum_{j=1}^p
\alpha_j a_j(\nuo)m_{2j}
\right\}.
\]
Thus
\[
q(\omega)
=
\GIG(\lambda,\chi,1),
\]
with
\begin{align}
\lambda
&=
\frac{\nuo+3}{2},
\\
\chi
&=
N
\sum_{j=1}^p
\alpha_j a_j(\nuo)m_{2j}.
\end{align}
The shared inverse moment is
\[
\eta
=
\frac{1}{\sqrt{\chi}}
\frac{
K_{(\nuo+1)/2}(\sqrt{\chi})
}{
K_{(\nuo+3)/2}(\sqrt{\chi})
}.
\]
If \(\chi\) is numerically close to zero, use
\[
\eta=\frac{1}{\nuo+1}.
\]

The posterior inclusion probabilities are updated using
\[
\alpha_j
=
\frac{
\pi_j\exp(\mathcal G_j)
}{
\sum_{k=1}^p
\pi_k\exp(\mathcal G_k)
}.
\]

\subsection*{Shared-\(\omega\) VEB algorithm for \(L=1\)}

\begin{enumerate}[leftmargin=2em]

\item \textbf{Input:}
summary statistic \(v\), sample size \(N\), reference LD matrix \(\Rbar\), prior weights \(\pi_j\), initial values \(\nuo^{(0)}\) and \({\sigb}^{(0)}\).

\item \textbf{Initialize:}
\[
\alpha_j^{(0)}=\pi_j,
\qquad
\eta^{(0)}=\frac{1}{\nuo^{(0)}+1}.
\]

\item \textbf{Repeat until convergence:}

\begin{enumerate}[leftmargin=2em]

\item For each SNP \(j=1,\dots,p\), compute
\[
a_j(\nuo)=\nuo\Rbar_{jj}+Nv_j^2.
\]

\item Update \(q_j(\beta)\) using the shared \(\eta\):
\[
s_j^2
=
\left[
\frac{1}{\sigb}
+
N a_j(\nuo)\eta
\right]^{-1},
\qquad
\mu_j=s_j^2Nv_j,
\qquad
m_{2j}=\mu_j^2+s_j^2.
\]

\item Update the posterior inclusion probabilities:
\[
\mathcal G_j
=
\frac12
\left[
\log\left(\frac{s_j^2}{\sigb}\right)
+
1
-
\frac{m_{2j}}{\sigb}
\right]
+
Nv_j\mu_j
-
\frac N2a_j(\nuo)m_{2j}\eta,
\]
\[
\alpha_j
=
\frac{
\pi_j\exp(\mathcal G_j)
}{
\sum_{k=1}^p
\pi_k\exp(\mathcal G_k)
}.
\]

\item Update the shared \(q(\omega)\):
\[
\lambda=\frac{\nuo+3}{2},
\qquad
\chi
=
N
\sum_{j=1}^p
\alpha_j a_j(\nuo)m_{2j},
\qquad
q(\omega)=\GIG(\lambda,\chi,1).
\]

\item Update the shared inverse moment:
\[
\eta
=
\frac{1}{\sqrt{\chi}}
\frac{
K_{(\nuo+1)/2}(\sqrt{\chi})
}{
K_{(\nuo+3)/2}(\sqrt{\chi})
}.
\]
If \(\chi\) is numerically close to zero, use
\[
\eta=\frac{1}{\nuo+1}.
\]

\item Update the prior variance:
\[
\sigb
\leftarrow
\sum_{j=1}^p
\alpha_j m_{2j}.
\]

\item Update the LD-mismatch parameter:
\[
\nuo
\leftarrow
\arg\max_{\nu>0}
\mathcal L_{\mathrm{shared}}(\nu,\sigb,q).
\]

\end{enumerate}

\item \textbf{Output:}
\[
\widehat\nuo,
\qquad
\widehat{\sigb},
\qquad
\alpha_j,
\qquad
q_j(\beta)=\Normal(\mu_j,s_j^2),
\qquad
q(\omega)=\GIG(\lambda,\chi,1).
\]

\end{enumerate}

\section{Variant replacing \(\omega\) by \(\nuo+3\)}

An even simpler approximation replaces the random variable
\[
\omega\sim\chisq_{\nuo+3}
\]
by its mean,
\[
\E(\omega)=\nuo+3.
\]
Equivalently, replace every occurrence of \(\omega^{-1}\) in the quadratic penalty by
\[
\frac{1}{\nuo+3}.
\]

For \(\gamma=j\), this gives the plug-in approximation
\[
p(v\mid \gamma=j,\beta,\nuo)
\approx
p_0(v\mid\nuo)
\exp\left\{
N\beta v_j
-
\frac{N a_j(\nuo)\beta^2}{2(\nuo+3)}
\right\}.
\]
There is no variational distribution for \(\omega\) and no KL penalty for \(\omega\).

Use the variational family
\[
q(\gamma=j,\beta)
=
\alpha_j q_j(\beta),
\qquad
q_j(\beta)=\Normal(\mu_j,s_j^2).
\]
The plug-in objective is
\begin{align}
\mathcal L_{\mathrm{plug}}(\nuo,\sigb,q)
&=
\log p_0(v\mid\nuo)
+
\sum_{j=1}^p
\alpha_j
\left[
\log \pi_j
-
\log \alpha_j
+
\mathcal H_j
\right],
\end{align}
where
\begin{align}
\mathcal H_j
&=
\frac12
\left[
\log\left(\frac{s_j^2}{\sigb}\right)
+
1
-
\frac{m_{2j}}{\sigb}
\right]
+
Nv_j\mu_j
-
\frac{N a_j(\nuo)m_{2j}}{2(\nuo+3)}.
\end{align}

For fixed \(\nuo\) and \(\sigb\), the coordinate update for \(q_j(\beta)\) is
\begin{align}
s_j^2
&=
\left[
\frac{1}{\sigb}
+
\frac{N a_j(\nuo)}{\nuo+3}
\right]^{-1},
\\
\mu_j
&=
s_j^2Nv_j,
\\
m_{2j}
&=
\mu_j^2+s_j^2.
\end{align}
The posterior inclusion probabilities are
\[
\alpha_j
=
\frac{
\pi_j\exp(\mathcal H_j)
}{
\sum_{k=1}^p
\pi_k\exp(\mathcal H_k)
}.
\]

\subsection*{Plug-in-\(\omega\) VEB algorithm for \(L=1\)}

\begin{enumerate}[leftmargin=2em]

\item \textbf{Input:}
summary statistic \(v\), sample size \(N\), reference LD matrix \(\Rbar\), prior weights \(\pi_j\), initial values \(\nuo^{(0)}\) and \({\sigb}^{(0)}\).

\item \textbf{Initialize:}
\[
\alpha_j^{(0)}=\pi_j.
\]

\item \textbf{Repeat until convergence:}

\begin{enumerate}[leftmargin=2em]

\item For each SNP \(j=1,\dots,p\), compute
\[
a_j(\nuo)=\nuo\Rbar_{jj}+Nv_j^2.
\]

\item Update \(q_j(\beta)\):
\[
s_j^2
=
\left[
\frac{1}{\sigb}
+
\frac{N a_j(\nuo)}{\nuo+3}
\right]^{-1},
\qquad
\mu_j=s_j^2Nv_j,
\qquad
m_{2j}=\mu_j^2+s_j^2.
\]

\item Compute the plug-in local objective:
\[
\mathcal H_j
=
\frac12
\left[
\log\left(\frac{s_j^2}{\sigb}\right)
+
1
-
\frac{m_{2j}}{\sigb}
\right]
+
Nv_j\mu_j
-
\frac{N a_j(\nuo)m_{2j}}{2(\nuo+3)}.
\]

\item Update posterior inclusion probabilities:
\[
\alpha_j
=
\frac{
\pi_j\exp(\mathcal H_j)
}{
\sum_{k=1}^p
\pi_k\exp(\mathcal H_k)
}.
\]

\item Update the prior variance:
\[
\sigb
\leftarrow
\sum_{j=1}^p
\alpha_j m_{2j}.
\]

\item Update the LD-mismatch parameter:
\[
\nuo
\leftarrow
\arg\max_{\nu>0}
\mathcal L_{\mathrm{plug}}(\nu,\sigb,q).
\]

\end{enumerate}

\item \textbf{Output:}
\[
\widehat\nuo,
\qquad
\widehat{\sigb},
\qquad
\alpha_j,
\qquad
q_j(\beta)=\Normal(\mu_j,s_j^2).
\]

\end{enumerate}

\section{SuSiE version with a shared \(q(\omega)\)}

Now consider a multi-effect model with
\[
b=\sum_{\ell=1}^L b_\ell,
\qquad
b_\ell=\beta_\ell e_{\gamma_\ell}.
\]
For effect \(\ell\), let
\[
P(\gamma_\ell=j)=\pi_{\ell j},
\qquad
\beta_\ell\mid\gamma_\ell=j,\sigma_\ell^2
\sim
\Normal(0,\sigma_\ell^2).
\]

Write
\[
A_\nu=\nuo\Rbar+Nvv^\top,
\qquad
Q_\nu(b)=b^\top A_\nu b.
\]
The collapsed lower bound with a single shared \(q(\omega)\) is
\begin{align}
\log p(v\mid b,\nuo)
&\geq
\log p_0(v\mid\nuo)
+
Nb^\top v
-
\frac N2\eta b^\top A_\nu b
-
\KL\{q(\omega)\|\chisq_{\nuo+3}\},
\end{align}
where
\[
\eta=\E_q(\omega^{-1}).
\]

Use the mean-field family
\[
q(b,\omega)
=
q(\omega)
\prod_{\ell=1}^L
q_\ell(\gamma_\ell,\beta_\ell),
\]
with
\[
q_\ell(\gamma_\ell=j,\beta_\ell)
=
\alpha_{\ell j}q_{\ell j}(\beta_\ell),
\qquad
q_{\ell j}(\beta_\ell)
=
\Normal(\mu_{\ell j},s_{\ell j}^2).
\]
Define
\[
m_{\ell j}=\alpha_{\ell j}\mu_{\ell j},
\qquad
m_\ell=(m_{\ell 1},\dots,m_{\ell p})^\top,
\qquad
m=\sum_{\ell=1}^L m_\ell,
\]
and
\[
m_{2,\ell j}
=
\E_{q_{\ell j}}(\beta_\ell^2)
=
\mu_{\ell j}^2+s_{\ell j}^2.
\]

The shared-\(\omega\) SuSiE ELBO can be written as
\begin{align}
\mathcal L_{\mathrm{SuSiE}}(\nuo,\{\sigma_\ell^2\},q)
&=
\log p_0(v\mid\nuo)
-
\KL\{q(\omega)\|\chisq_{\nuo+3}\}
\nonumber\\
&\quad
+
N v^\top m
-
\frac N2\eta S_Q
\nonumber\\
&\quad
+
\sum_{\ell=1}^L
\sum_{j=1}^p
\alpha_{\ell j}
\left[
\log\pi_{\ell j}
-
\log\alpha_{\ell j}
\right.
\nonumber\\
&\qquad\qquad\qquad\left.
+
\frac12
\left\{
\log\left(\frac{s_{\ell j}^2}{\sigma_\ell^2}\right)
+
1
-
\frac{m_{2,\ell j}}{\sigma_\ell^2}
\right\}
\right],
\label{eq:susie_shared_elbo}
\end{align}
where
\[
S_Q
=
\E_q\{b^\top A_\nu b\}.
\]
An efficient expression for this expectation is
\[
S_Q
=
m^\top A_\nu m
+
\sum_{\ell=1}^L
\left[
\sum_{j=1}^p
\alpha_{\ell j}m_{2,\ell j}(A_\nu)_{jj}
-
m_\ell^\top A_\nu m_\ell
\right].
\]
The second term adds back the within-effect posterior variance and
removes the squared posterior mean already counted in \(m^\top A_\nu m\).

\subsection*{Coordinate update for one effect}

When updating effect \(\ell\), hold the other effects and \(q(\omega)\)
fixed. Let
\[
m_{-\ell}
=
\sum_{k\neq\ell}m_k
=
m-m_\ell,
\qquad
c_\ell=A_\nu m_{-\ell}.
\]
For candidate SNP \(j\), the terms involving \(\beta_\ell\) are
\[
N\beta_\ell v_j
-
N\eta\beta_\ell (c_\ell)_j
-
\frac N2\eta (A_\nu)_{jj}\beta_\ell^2
+
\log\Normal(\beta_\ell;0,\sigma_\ell^2).
\]
Thus
\[
q_{\ell j}(\beta_\ell)
=
\Normal(\mu_{\ell j},s_{\ell j}^2),
\]
where
\begin{align}
s_{\ell j}^2
&=
\left[
\frac{1}{\sigma_\ell^2}
+
N\eta (A_\nu)_{jj}
\right]^{-1},
\\
\mu_{\ell j}
&=
s_{\ell j}^2
N
\left[
v_j-\eta(c_\ell)_j
\right],
\\
m_{2,\ell j}
&=
\mu_{\ell j}^2+s_{\ell j}^2.
\end{align}

The coordinate score for \(\gamma_\ell=j\) is
\begin{align}
\mathcal D_{\ell j}
&=
\frac12
\left[
\log\left(\frac{s_{\ell j}^2}{\sigma_\ell^2}\right)
+
1
-
\frac{m_{2,\ell j}}{\sigma_\ell^2}
\right]
\nonumber\\
&\quad
+
N\mu_{\ell j}
\left[
v_j-\eta(c_\ell)_j
\right]
-
\frac N2\eta (A_\nu)_{jj}m_{2,\ell j}.
\end{align}
Therefore,
\[
\alpha_{\ell j}
=
\frac{
\pi_{\ell j}\exp(\mathcal D_{\ell j})
}{
\sum_{k=1}^p
\pi_{\ell k}\exp(\mathcal D_{\ell k})
}.
\]

\subsection*{Shared \(q(\omega)\) update}

For fixed effect factors, the optimal shared \(q(\omega)\) is
\[
q(\omega)
\propto
\chisq_{\nuo+3}(\omega)
\exp\left\{
-
\frac{N S_Q}{2\omega}
\right\}.
\]
Thus
\[
q(\omega)=\GIG(\lambda,\chi,1),
\]
with
\begin{align}
\lambda
&=
\frac{\nuo+3}{2},
\\
\chi
&=
N S_Q.
\end{align}
The shared inverse moment is
\[
\eta
=
\frac{1}{\sqrt{\chi}}
\frac{
K_{(\nuo+1)/2}(\sqrt{\chi})
}{
K_{(\nuo+3)/2}(\sqrt{\chi})
}.
\]
If \(\chi\) is numerically close to zero, use
\[
\eta=\frac{1}{\nuo+1}.
\]

\subsection*{Empirical-Bayes updates}

For fixed \(\nuo\), a natural update for the effect-specific prior
variance is
\[
\sigma_\ell^2
\leftarrow
\sum_{j=1}^p
\alpha_{\ell j}m_{2,\ell j}.
\]
There is generally no closed-form update for \(\nuo\) in this
shared-\(\omega\) collapsed objective. A practical update is
\[
\nuo
\leftarrow
\arg\max_{\nu>0}
\mathcal L_{\mathrm{SuSiE}}(\nu,\{\sigma_\ell^2\},q),
\]
or a grid search over candidate values of \(\nu\). When \(\nuo\) changes,
\(A_\nu\), \(\lambda\), \(\chi\), \(\eta\), and the null term
\(\log p_0(v\mid\nu)\) should be recomputed.

\subsection*{Shared-\(\omega\) collapsed SuSiE algorithm}

\begin{enumerate}[leftmargin=2em]

\item \textbf{Input:}
summary statistic \(v\), sample size \(N\), reference LD matrix \(\Rbar\),
number of effects \(L\), prior weights \(\pi_{\ell j}\), initial values
\(\nuo^{(0)}\) and \((\sigma_\ell^2)^{(0)}\).

\item \textbf{Initialize:}
for each effect \(\ell\), initialize
\[
\alpha_{\ell j}^{(0)}=\pi_{\ell j},
\qquad
\mu_{\ell j}^{(0)}=0,
\qquad
s_{\ell j}^{2(0)}=(\sigma_\ell^2)^{(0)}.
\]
Set
\[
\eta^{(0)}=\frac{1}{\nuo^{(0)}+1}.
\]

\item \textbf{Repeat until convergence:}

\begin{enumerate}[leftmargin=2em]

\item Form
\[
A_\nu=\nuo\Rbar+Nvv^\top.
\]

\item For each effect \(\ell=1,\dots,L\), perform an IBSS update:

\begin{enumerate}[leftmargin=2em]

\item Compute
\[
m_\ell=(\alpha_{\ell 1}\mu_{\ell 1},\dots,\alpha_{\ell p}\mu_{\ell p})^\top,
\qquad
m_{-\ell}=\sum_{k\neq\ell}m_k,
\qquad
c_\ell=A_\nu m_{-\ell}.
\]

\item For each SNP \(j=1,\dots,p\), update
\[
s_{\ell j}^2
=
\left[
\frac{1}{\sigma_\ell^2}
+
N\eta (A_\nu)_{jj}
\right]^{-1},
\]
\[
\mu_{\ell j}
=
s_{\ell j}^2
N
\left[
v_j-\eta(c_\ell)_j
\right],
\]
\[
m_{2,\ell j}
=
\mu_{\ell j}^2+s_{\ell j}^2.
\]

\item Compute \(\mathcal D_{\ell j}\) and update
\[
\alpha_{\ell j}
=
\frac{
\pi_{\ell j}\exp(\mathcal D_{\ell j})
}{
\sum_{k=1}^p
\pi_{\ell k}\exp(\mathcal D_{\ell k})
}.
\]

\end{enumerate}

\item Recompute
\[
m=\sum_{\ell=1}^L m_\ell
\]
and
\[
S_Q
=
m^\top A_\nu m
+
\sum_{\ell=1}^L
\left[
\sum_{j=1}^p
\alpha_{\ell j}m_{2,\ell j}(A_\nu)_{jj}
-
m_\ell^\top A_\nu m_\ell
\right].
\]

\item Update the shared \(q(\omega)\):
\[
\lambda=\frac{\nuo+3}{2},
\qquad
\chi=N S_Q,
\qquad
q(\omega)=\GIG(\lambda,\chi,1).
\]
Then update
\[
\eta=\E_q(\omega^{-1}).
\]

\item Optionally update
\[
\sigma_\ell^2
\leftarrow
\sum_{j=1}^p
\alpha_{\ell j}m_{2,\ell j},
\qquad
\ell=1,\dots,L.
\]

\item Optionally update \(\nuo\) by grid search or one-dimensional
optimization of \(\mathcal L_{\mathrm{SuSiE}}\).

\end{enumerate}

\item \textbf{Output:}
\[
\widehat\nuo,
\qquad
\widehat\sigma_1^2,\dots,\widehat\sigma_L^2,
\qquad
\alpha_{\ell j},
\qquad
q_{\ell j}(\beta_\ell)=\Normal(\mu_{\ell j},s_{\ell j}^2),
\qquad
q(\omega)=\GIG(\lambda,\chi,1).
\]

The posterior inclusion probability for SNP \(j\) is
\[
\mathrm{PIP}_j
=
1-
\prod_{\ell=1}^L
(1-\alpha_{\ell j}).
\]

\end{enumerate}

\section{Comments}

For \(L=1\), exact one-dimensional integration over \(\beta\) is also feasible. The VEB approximation is still useful because it provides a computational template that generalizes naturally to multi-effect SuSiE-style models.

The empirical-Bayes estimate \(\widehat\nuo\) should be interpreted as a working LD-mismatch parameter. Smaller values of \(\nuo\) allow more uncertainty around the reference LD matrix \(\Rbar\), whereas larger values concentrate \(R\) more tightly around \(\Rbar\).

\end{document}
